\section{適用実験}

\subsection{データセット}
評価には，2つのREST APIを対象とした計26ケースを用いた．1ケースは，OSSリポジトリから収集した誤用コードと，そのリポジトリ上で実際に行われた修正後のコードの組であり，PythonやTypeScriptなど8言語にまたがる．26ケースからは，エンドポイントの誤用とリクエストヘッダの誤用が計31件得られた．内訳を\textbf{表\ref{tab:dataset}}に示す．

SwitchBot APIは，v1.0からv1.1への移行を対象とする．この移行では，エンドポイントのバージョンが更新されるとともに，トークンのみの認証方式が署名に基づく認証方式へ変更される．8ケースはいずれも同一の修正を要し，各ケースが両方の誤用を含む．

Fitbit APIは2種類の移行を含む．1つはFitbit API内のバージョン移行であり，Sleep APIのv1からv1.2への更新やOAuth 1.0からOAuth 2.0への変更が対象である．もう1つは，Fitbit Web APIの廃止に伴うGoogle Health APIへの移行であり，移行先が別サービスのAPIにまたがる．このデータセットを準備するための作業は全て著者が行った．

\begin{table}[t]
\centering
\caption{データセットに含まれる誤用の件数}
\label{tab:dataset}
\small
\begin{tabular}{l|rr}
\hline\hline
 & SwitchBot API & Fitbit API \\
\hline
エンドポイント & 8 & 11 \\
リクエストヘッダ & 8 & 4 \\
\hline
合計 & 16 & 15 \\
\hline
\end{tabular}
\end{table}

\subsection{比較手法}

BL-Webは，仕様のリンクをプロンプトに与え，そのリンク先の本文を取得して使用される．本文をあらかじめプロンプトへ埋め込むのではなく，リンクを起点に取得させる点が提案手法と異なる．
DsLsは先行研究\cite{inoue2025dbrag}で最良であった構成であり，非推奨仕様と最新仕様をそれぞれ自然言語テキストとコード片に分けた4つのデータベースを個別に検索し，その結果を連結して与える．
なお，Fitbitでは，API内の移行とサービス間の移行で参照すべき仕様が異なるため，DsLsの索引を移行の種類ごとに分けた．


\subsection{評価方法と実験環境}
LLMの出力の揺れを抑えるため，手法とケースの組ごとに5回生成した．

修正の成否は，誤用コードと修正後のコードの差分から，除去されるべき記述と導入されるべき記述を修正項目として導出し，生成されたパッチがそのすべてを満たすかによって判定した．手法ごとに，全項目を満たした試行の割合を修正率とする．

生成には研究室内に設置したLM Studio上の\texttt{gpt-oss-120b}を用いた．埋め込みには768次元の\texttt{nomic-embed-text-v1.5}を，グラフストアにはNeo4jを用いた．すべての生成をローカル環境で実行しており，外部のLLMサービスは利用していない．
